\documentclass[../../../../hsm_tok_q.tex]{subfiles}

\begin{document}

    \textsf{図 \ref{fig/hsm_tok_2014_1st_01_09_01_q}}で，四角形ABCDは平行四辺形である．

    解答欄に示した図をもとにして，辺AD上にあり，
    頂点B，頂点Cまでの距離が等しい点Pを，
    定規とコンパスを用いて作図によって求め，
    点Pの位置を示す文字Pも書け．
    
    ただし，作図に用いた線は消さないでおくこと．
    %
    \vspace{.5\baselineskip}
    {\par \centering
        \includegraphics%
            {fig_hsm_tok_2014_1st_01_09/fig_hsm_tok_2014_1st_01_09_01_q.pdf}
            \captionof{figure}{} \label{fig/hsm_tok_2014_1st_01_09_01_q}
    \par}
    %
    \vspace{.2\baselineskip}
\end{document}
